% ======================================================================
% Chapter 4 -- Part D: business case and rollout; changed requirement;
% end matter
% ======================================================================

\section{Business Case: Is It Worth It?}
\label{sec:policy-business}

\subsection{The value case}

The value case uses the method of Chapter~\ref{ch:business}, with every input marked measured or assumed in \texttt{ch04/value\_ch04.py}. Its base result (Table~\ref{tab:policy-value}) is modest, and it is worth understanding why.

\begin{thinkabout}
The floor case pays back in about 18 months. For an HR tool, is that good or bad? What else would you want to know before answering?
\end{thinkabout}

\begin{table}[htbp]
\centering
\small
\caption{Annual value case for the grounded design (base assumptions). One-time build: \$120{,}000.}
\label{tab:policy-value}
\begin{tabularx}{\textwidth}{@{}>{\raggedright\arraybackslash}p{3.6cm}r>{\raggedright\arraybackslash}X@{}}
\toprule
\rowcolor{coHead} Line & \$ per year & Main input \\
\midrule
HR capacity (60\% redeployed) & 97{,}800 & 50\% resolution rate (\emph{assumed}); 11 min per ticket (measured) \\
Fewer wrong answers & 74{,}900 & 7\% $\to$ 3\% on resolved questions; \$120 per error (estimated) \\
Run cost & $-94{,}000$ & 48{,}000 questions at \$0.02; platform; 0.3 engineer; 0.2 HR curator \\
\midrule
\rowcolor{coGreenBg} Net (floor case) & 78{,}700 & Payback about 18 months \\
\bottomrule
\end{tabularx}
\end{table}

Three features of the case are typical of information products:
\begin{itemize}[itemsep=2pt]
  \item \emph{Most of the cost is people, not models.} Of the \$94{,}000 run cost, under \$1{,}000 is search and model calls. The rest is the engineer who keeps the system healthy and the HR curator who keeps metadata right. A team that budgets only for compute will be surprised.
  \item \emph{The case swings on one assumed rate.} At a 30\% resolution rate the net value is about \$9{,}700 a year; at 70\% it is about \$148{,}000. The pilot must measure this rate honestly, using the seven-day rule.
  \item \emph{The largest benefit is not in the floor.} Employees spend about 12 minutes searching and filing each ticket before any waiting begins. If half of those tickets are no longer needed, about \$225{,}000 a year of engineers' time is freed---more than the whole floor case. It is left out of the floor because it is diffuse and hard to redeploy, but it belongs in the memo as upside, along with the reduced legal exposure from wrong answers about pay and leave.
\end{itemize}

A sponsor deciding on an 18-month payback with significant upside and a strong risk argument has a reasonable decision to make. A sponsor shown only the upside would have been misled.

\subsection{Pilot plan}

\begin{itemize}[itemsep=1pt]
  \item \emph{Scope:} the UK office (700 employees), six weeks; all policies, all grades.
  \item \emph{Comparison:} policy tickets per 100 employees in the US and German offices over the same weeks, and in the UK office over the six weeks before launch.
  \item \emph{Owner:} Director of People Operations.
  \item \emph{Stop rules:} any confirmed leakage of a restricted document; any answer applying another country's pay or leave rules that reaches an employee and is acted on.
\end{itemize}

\begin{table}[htbp]
\centering
\small
\caption{Go/no-go criteria for the UK pilot.}
\label{tab:policy-gonogo}
\begin{tabularx}{\textwidth}{@{}>{\raggedright\arraybackslash}X>{\raggedright\arraybackslash}p{2.4cm}>{\raggedright\arraybackslash}p{3.2cm}@{}}
\toprule
\rowcolor{coHead} Criterion & Baseline & Go if \\
\midrule
\kpi{Resolved without HR} UK policy tickets per 100 employees, vs.\ comparison offices & --- & $\ge$ 40\% lower \\
\kpi{Answer accuracy} Answer accuracy, weekly sampled review & 93\% (HR) & $\ge$ 97\% \\
\kpi{Time to answer} Median time to answer, questions passed to HR & 1.8 days & $\le$ 1 day \\
\risk{Restricted leak} Leakage (evaluation and production) & --- & 0 (stop rule) \\
Abstention recall on evaluation set & --- & $\ge$ 0.90 \\
Share of UK employees who used it at least twice & --- & $\ge$ 40\% \\
\bottomrule
\end{tabularx}
\end{table}

\subsection{The memo}

\begin{tcolorbox}[colback=white, colframe=black!60, boxrule=0.6pt, arc=1pt, left=10pt, right=10pt, top=8pt, bottom=8pt, fontupper=\small]
\textbf{To:} Director, People Operations \hfill \textbf{From:} Policy Assistant Team\\
\textbf{Re:} Policy answers in a minute, for your own country and grade --- request to pilot

\smallskip\hrule\smallskip

\textbf{Recommendation.} Approve a six-week pilot in the UK office; build cost \$120{,}000; run cost about \$94{,}000 a year, mostly staff.

\textbf{The problem.} 2{,}600 policy tickets a month wait a median of 1.8 days. 7\% of answers are wrong, mostly because the wrong country's or an outdated policy was applied. Generalists spend most of their week on repeat questions.

\textbf{What changes.} Employees ask in the intranet and get an answer drawn only from the policies that apply to them, with a link to the exact section and version. When the policies do not answer the question, the assistant says so and offers to send it to HR. Manager-only documents are never searched for other employees. A published policy is answerable within an hour.

\textbf{Value.} \emph{Floor case:} about \$79{,}000 a year from HR time and fewer wrong answers; payback about 18 months. \emph{Upside:} about \$225{,}000 a year of employees' time no longer spent searching and filing tickets; lower exposure from wrong pay and leave answers.

\textbf{Risks and safeguards.} Exposure of restricted documents is prevented by filtering before search, tested on every release; any occurrence stops the pilot. Logs are kept 30 days and reported in aggregate only, as agreed with the works council's representatives.

\textbf{Decision at six weeks.} Go if UK tickets fall at least 40\% against other offices, sampled accuracy is at least 97\%, and leakage is zero. The pilot will also show whether curated answers for the most common questions are worth adding.

\textbf{What we need from you.} The policy team's time to tag 340 documents (about two weeks); agreement that generalists' freed time goes to manager support.
\end{tcolorbox}

% ----------------------------------------------------------------------
\section{Changed Requirement: What If the World Changes?}
\label{sec:policy-changed}

Choose one of the following changes. Redo the design table, the trade-off table, the failure walk-through and the memo. State which earlier decisions survive unchanged and which do not.

\begin{enumerate}[itemsep=3pt]
  \item \textbf{A new constraint from the works council.} The German works council requires that no question text from German employees be stored at all, even for 30 days, and that no evaluation use real German questions. How do you keep measuring \kpi{Answer accuracy} for German employees? Does this change the deployment, the evaluation, or the design itself?
  \item \textbf{Ten times the users.} Alder \& Finch acquires a firm with 28{,}000 employees in nine more countries. Policies now differ across twelve countries and four collective agreements, and the library grows to 2{,}500 documents. Which parts of the grounded design scale unchanged? Does the case for the curated design's curated layer get stronger or weaker? What happens to the metadata burden?
  \item \textbf{Half the budget.} The sponsor approves the pilot but only \$50{,}000 a year to run it. The engineer and curator lines dominate the run cost. What do you cut, what risk does each cut raise, and what would you tell the sponsor?
\end{enumerate}

% ----------------------------------------------------------------------
\section*{Chapter Summary}
\addcontentsline{toc}{section}{Chapter Summary}

\begin{itemize}[itemsep=2pt]
  \item The \texttt{rag-agent-search} sample is a sound starting point that leaves six decisions open: grounding, citations, user identity, access control, freshness and evaluation.
  \item For information products, deflection is a trap unless it counts only resolved questions and is paired with sampled accuracy.
  \item Access control in retrieval belongs before search, in code, driven by verified identity and document metadata. Metadata becomes part of the system's correctness and needs an owner and validation.
  \item When designs tie within the noise and one extends the other, build the smaller first and let the pilot's data decide on the extension.
  \item A first deployment should settle residency, least privilege, pinned model versions, private access and event-driven freshness.
  \item Measure retrieval separately from answers. Leakage is measured at retrieval, because a restricted passage that reaches the model has already crossed its boundary.
  \item Information products often have modest hard-dollar cases whose costs are mostly people; lead with the floor and present diffuse time savings and risk reduction as upside.
\end{itemize}

\section*{Review Questions}
\addcontentsline{toc}{section}{Review Questions}

\begin{enumerate}[itemsep=2pt]
  \item List the six open decisions in the sample. Which one removes the sample design from consideration, and why is it a constraint rather than a criterion?
  \item Why is the seven-day rule in the definition of \kpi{Resolved without HR} necessary? What failure would deflection alone hide?
  \item Why does the grounded design filter before retrieval rather than check the answer after generation?
  \item In Listing~\ref{lst:search-tool}, what would go wrong if the employee's country were a parameter of \texttt{search\_policies} chosen by the model?
  \item The grounded design and the curated design nearly tie. Give the three questions used to settle the tie, and the answer to each.
  \item What does the citation guardrail check, and what does it not check? How does the decision record handle that gap?
  \item Explain why the evaluation measures leakage at retrieval. Use the demonstration output in your answer.
  \item Why is the model name pinned to a specific version in deployment?
  \item Most of the run cost is staff. What does that imply for how the team should present the cost to the sponsor?
  \item Which stakeholder introduced a constraint on evaluation, and how did the design respond?
\end{enumerate}

\section*{Exercises}
\addcontentsline{toc}{section}{Exercises}

\paragraph{Design track}
\begin{enumerate}[itemsep=2pt]
  \item \textbf{KPI tree.} Draw the KPI tree described in Section~\ref{sec:policy-kpis}, with leaves taken from Table~\ref{tab:policy-thresholds}.
  \item \textbf{Effective dates.} Policies are often published before they take effect. Extend the grounded design so that employees get the policy in force today by default but can ask about an upcoming change. Which steps change?
  \item \textbf{Component cards.} Write component cards for \texttt{search\_policies} and the citation check.
  \item \textbf{Changed requirement.} Complete one of the three changes in Section~\ref{sec:policy-changed}.
\end{enumerate}

\paragraph{Build track}
\begin{enumerate}[resume, itemsep=2pt]
  \item \textbf{Fail closed.} Write unit tests for \texttt{build\_filter} and \texttt{search\_policies}, including a test that proves no search runs without a profile.
  \item \textbf{Claim-level grounding.} Implement a check that asks a second model whether each cited passage supports the sentence that cites it. Measure its agreement with human reviewers and its added latency and cost on the evaluation set, and recommend whether to adopt it.
  \item \textbf{Freshness probe.} Implement the hourly freshness probe from the failure walk-through and alert when a newly published version is not retrievable within an hour.
\end{enumerate}

% ----------------------------------------------------------------------
\section*{Lab 4: Deploy and Measure the Policy Assistant}
\addcontentsline{toc}{section}{Lab 4: Deploy and Measure the Policy Assistant}

\paragraph{Goal.} Deploy the sample design and the grounded design against the same synthetic policy library, run the evaluation set against both, and explain the difference in terms of design decisions.

\paragraph{Time.} About three hours, including about 30 minutes waiting for infrastructure.

\paragraph{Steps.}
\begin{enumerate}[itemsep=2pt]
  \item In \texttt{ch04/}, run the setup script. It copies the pinned sample, creates a Cloud Storage bucket, and loads a synthetic Alder \& Finch library of 60 policies for three countries, with manifests for both the unstructured schema of the sample design and the structured schema of the grounded design.
  \item Provision two search data stores with the sample's Terraform, one per schema. Note the time each takes to become queryable.
  \item Run the sample design locally with \texttt{adk web} and ask three questions from the evaluation set as a UK, a US and a German employee would. Record what changes and what does not.
  \item Build the grounded design from the provided skeleton by completing \texttt{build\_filter}, \texttt{search\_policies} and \texttt{citation\_check}. Run the unit tests.
  \item Deploy the grounded design to Cloud Run with the command in Section~\ref{sec:policy-deploy}. Confirm that an unauthenticated request is refused.
  \item Run \texttt{python run\_eval.py --design sample} and \texttt{--design grounded}. The script replays the 240 questions with synthetic employee profiles and writes evaluation records; \texttt{rag\_metrics.py} turns them into Table~\ref{tab:policy-thresholds}'s metrics.
  \item Publish a changed policy to the bucket and measure how long until the grounded design answers with the new version.
  \item Run the cleanup script. Check in the console that the data stores, bucket and Cloud Run services are gone.
\end{enumerate}

\paragraph{Deliverable.} A results table for the sample design and the grounded design against every threshold; for each failed threshold, the step responsible and the change you would make; and a one-paragraph answer to: \emph{should the curated layer of the curated design be built?}

% ----------------------------------------------------------------------
\section*{Further Reading}
\addcontentsline{toc}{section}{Further Reading}

\begin{itemize}[itemsep=2pt]
  \item \citet{lewis2020rag} remains the clearest statement of why retrieval and generation should be treated as separate components.
  \item \citet{es2024ragas} proposes reference-free metrics for RAG systems, including faithfulness and context relevance; compare them with the reference-based metrics used in this chapter.
  \item Agent Search's documentation on document schemas and filter expressions is the place to check which metadata fields can be filtered and with which operators; both change more often than the design principles in this chapter.
\end{itemize}
